\documentclass{article} % This command is used to set the type of
% document you are working on such as an article, book, or presenation

\usepackage{geometry} % This package allows the editing of the page layout
\usepackage{amsmath}  % This package allows the use of a large range
% of mathematical formula, commands, and symbols
\usepackage{graphicx}  % This package allows the importing of images
\usepackage{
  amsmath,      % Math Environments
  amssymb,      % Extended Symbols
  enumerate,        % Enumerate Environments
  graphicx,      % Include Images
  lastpage,      % Reference Lastpage
  multicol,      % Use Multi-columns
  multirow      % Use Multi-rows
}

\newcommand{\question}[2][]{
  \begin{flushleft}
    \textbf{Question #1}: \textit{#2}

\end{flushleft}}
\newcommand{\sol}{\textbf{Solution}:} %Use if you want a boldface solution line
\newcommand{\maketitletwo}[1][]{
  \begin{center}
    \Large{\textbf{Assignment #1}

    MATH4362} % Name of course here
    \vspace{5pt}

    \normalsize{N. Ohayon Rozanes % Your name here

    \today}        % Change to due date if preferred
    \vspace{15pt}

\end{center}}
\begin{document}
\maketitletwo[5]  % Optional argument is assignment number
%Keep a blank space between maketitletwo and \question[1]

\question[1]{}
We solve $u_t = u_{xx}$ on the real line with initial data
$f(x) = \delta(x-1) + \delta(x+1)$.
Taking the Fourier transform of the equation gives
$\frac{\partial\hat{u}}{\partial t} = -k^2\hat{u}$,
with solution $\hat{u}(t,k) = e^{-k^2 t}\hat{f}(k)$. We compute:
\[
  \mathcal{F}[f] = \mathcal{F}[\delta(x-1)] + \mathcal{F}[\delta(x+1)]
  = \frac{e^{-ik} + e^{ik}}{\sqrt{2\pi}} = \sqrt{\tfrac{2}{\pi}}\cos k
\]
Therefore:
\[
  u(t,x) = \frac{1}{\sqrt{2\pi}}\int_{-\infty}^{\infty} e^{ikx-k^2t}
  \sqrt{\tfrac{2}{\pi}}\cos k\,dk
  = \frac{1}{\pi}\int_{-\infty}^{\infty} e^{ikx-k^2t}\cos k\,dk
\]
Writing $\cos k = \frac{e^{ik}+e^{-ik}}{2}$ splits this into two
Gaussian integrals.
Using $\int_{-\infty}^{\infty}e^{-tk^2+iak}dk =
\sqrt{\tfrac{\pi}{t}}e^{-a^2/(4t)}$
(completed square), with $a = x\pm 1$:
\[
  u(t,x) = \frac{1}{2\pi}\sqrt{\frac{\pi}{t}}\left[e^{-(x+1)^2/(4t)}
  + e^{-(x-1)^2/(4t)}\right]
  = \frac{1}{2\sqrt{\pi t}}\left[e^{-(x+1)^2/(4t)} + e^{-(x-1)^2/(4t)}\right]
\]

\question[2]{}
We solve $u_t = u_{xx}$ on $-\infty < x < \infty$, $t > 0$ with
$u(0,x) = \delta'(x-\xi)$. Taking the Fourier transform of the initial
data:
\[
  \hat{f}(k) =
  \frac{1}{\sqrt{2\pi}}\int_{-\infty}^{\infty}\delta'(x-\xi)e^{ikx}dx
  = \frac{-ik}{\sqrt{2\pi}}e^{ik\xi}
\]
Therefore $\hat{u}(t,k) = \frac{-ik}{\sqrt{2\pi}}e^{ik\xi-k^2t}$, and
applying the inverse Fourier transform:
\[
  u(t,x) = \frac{-i}{2\pi}\int_{-\infty}^{\infty} k\,e^{ik(\xi-x)-k^2t}\,dk
\]

\question[3]{}
We solve the following PDE on the rectangle $-1 < x < 1, 0 < y < 1$
\[
  \begin{cases}
    u_t = 2\Delta u\\
    u(t,-1, y) = u(t, 1, y) = 0 & \forall y \\
    u_n(t, x,0) = u_n(t,x, 1) = 0 \\
  \end{cases}
\]
With initial distribution:
\[
  u(0,x,y) =
  \begin{cases}
    -1 & x < 0 \\
    1 & x > 0
  \end{cases}
\]
Using the seperable anstaz for the above:
\[
  u = e^{-\lambda t}p(x)q(y)
\]
We calculate partials and substitute into the PDE:
\[
  -\lambda pq = 2(p''q + pq'')
\]
Which, after some algebra, yeilds the following second order ODEs:
\[
  p'' +  \frac{\mu}{2}p = 0
\]
\[
  q'' +  \frac{\lambda-\mu}{2}q = 0
\]

For some real number $\mu$. Without considering the boundary
conditions, we have that the possible solutions are:

\[
  p = \{A\cos\!\left(\sqrt{\tfrac{\mu}{2}}\,x\right),
  B\sin\!\left(\sqrt{\tfrac{\mu}{2}}\,x\right)\}
\]

\[
  q = \{C\cos\!\left(\sqrt{\tfrac{\lambda-\mu}{2}}\,y\right) ,
  D\sin\!\left(\sqrt{\tfrac{\lambda-\mu}{2}}\,y\right)\}
\]

Consider that the normal vector at every every point on the top long
edge is the vector $(0,1)$ and $(0,-1)$ at every point of the bottom
long edge: Then we can find the equations of the
insulated condition as follows

\[
  \nabla u \cdot (0,1)\big|_{y=1} = u_y(t,x,1) = e^{-\lambda
  t}p(x)q'(1) = 0 \implies q'(1) = 0
\]
\[
  \nabla u \cdot (0,-1)\big|_{y=0} = -u_y(t,x,0) = -e^{-\lambda
  t}p(x)q'(0) = 0 \implies q'(0) = 0
\]

Applying this to the possible solutions of $q$, we have that

\[
  q' =
  -C\sqrt{\tfrac{\lambda-\mu}{2}}\sin\!\left(\sqrt{\tfrac{\lambda-\mu}{2}}\,y\right)
  +
  D\sqrt{\tfrac{\lambda-\mu}{2}}\cos\!\left(\sqrt{\tfrac{\lambda-\mu}{2}}\,y\right)
\]

At $y=0$, $\sin(0)=0$ and $\cos(0)=1$, so:
\[
  q'(0) = D\sqrt{\tfrac{\lambda-\mu}{2}} = 0 \implies D = 0
\]

With $D=0$, at $y=1$:
\[
  q'(1) =
  -C\sqrt{\tfrac{\lambda-\mu}{2}}\sin\!\left(\sqrt{\tfrac{\lambda-\mu}{2}}\right)
  = 0
\]

For a non-trivial solution we need
$\sin\!\left(\sqrt{\tfrac{\lambda-\mu}{2}}\right) = 0$, so
$\sqrt{\tfrac{\lambda-\mu}{2}} = n\pi$ for $n = 0, 1, 2, \ldots$,
giving eigenfunctions:
\[
  q_n(y) = \cos(n\pi y)
\]
We now condition the Dirichlet conditions on the short edges, which we get
\begin{align*}
  p(1) = A\cos\!\left(\sqrt{\frac{\mu}{2}}\right) +
  B\sin\!\left(\sqrt{\frac{\mu}{2}}\right) &=0 \\
  p(-1) = A\cos\!\left(\sqrt{\frac{\mu}{2}}\right)
  -B\sin\!\left(\sqrt{\frac{\mu}{2}}\right) &=0 \\
\end{align*}
Adding one to the other we have that
\[
  2A\cos\!\left(\sqrt{\frac{\mu}{2}}\right) = 0
\]
So either $A$ is trivial or $\sqrt{\frac{\mu}{2}} = \frac{\pi}{2} + n\pi$
for $n = 0,1,\hdots$

Substracting we get that
\[
  2B\sin\!\left(\sqrt{\frac{\mu}{2}}\right) =0
\]
So either $B$ is trivial or $\sqrt{\frac{\mu}{2}} = m\pi$ for $m =
1,2,\hdots$.

Recalling that $\sqrt{\tfrac{\lambda-\mu}{2}} = n\pi$, i.e.\ $\lambda
= \mu + 2n^2\pi^2$,
we can solve for $\lambda$ in each family. For the cosine family,
$\sqrt{\tfrac{\mu}{2}} = \tfrac{(2m-1)\pi}{2}$ gives $\mu =
\tfrac{(2m-1)^2\pi^2}{2}$, so
\[
  \lambda_{mn}^{(c)} = \frac{(2m-1)^2\pi^2}{2} + 2n^2\pi^2
\]
For the sine family, $\sqrt{\tfrac{\mu}{2}} = m\pi$ gives $\mu = 2m^2\pi^2$, so
\[
  \lambda_{mn}^{(s)} = 2m^2\pi^2 + 2n^2\pi^2 = 2\pi^2(m^2+n^2)
\]
Without accounting for the initial data, we get the general solution
\[
  u(t,x,y) = \sum_{m=1}^{\infty}\sum_{n=0}^{\infty}
  e^{-\left(\frac{(2m-1)^2\pi^2}{2} + 2n^2\pi^2\right)t}
  A_{mn}\cos\!\left(\tfrac{(2m-1)\pi}{2}x\right)\cos(n\pi y)
  +
  \sum_{m=1}^{\infty}\sum_{n=0}^{\infty}
  e^{-2\pi^2(m^2+n^2)t}
  B_{mn}\sin(m\pi x)\cos(n\pi y)
\]

Before solving for the coefficients using the initial data, we
formulate the function $f$ in a more workable manner using $\sigma$
functions:

\[
  f = -1 + 2\sigma(x)
\]

We are now ready to use the initial data to solve for the
coefficients. Setting $t=0$:
\[
  f(x) = \sum_{m=1}^{\infty}\sum_{n=0}^{\infty}
  A_{mn}\cos\!\left(\tfrac{(2m-1)\pi}{2}x\right)\cos(n\pi y)
  +
  \sum_{m=1}^{\infty}\sum_{n=0}^{\infty}
  B_{mn}\sin(m\pi x)\cos(n\pi y)
\]
Since $f$ is independent of $y$, multiplying both sides by $\cos(k\pi
y)$ and integrating over $[0,1]$ gives zero on the left for any $k\neq
0$, so orthogonality of $\{\cos(n\pi y)\}$ forces $A_{mn}=B_{mn}=0$
for all $n\neq 0$. Furthermore, $f$ is odd in $x$ while
$\cos\!\left(\tfrac{(2m-1)\pi}{2}x\right)$ is even, so orthogonality
in $x$ forces $A_{m0}=0$ for all $m$. The series collapses to:
\[
  f(x) = \sum_{m=1}^{\infty} B_m \sin(m\pi x)
\]

This is the ordinary sine Fourier series of $f$, with coefficients:
\[
  B_m = \int_{-1}^1 f(x)\sin(m\pi x)\,dx
\]
Since $f(x)\sin(m\pi x)$ is even, we have that
\[
  B_m = 2\int_0^1 \sin(m\pi x)\,dx
  = \left[-\frac{2\cos(m\pi x)}{m\pi}\right]_0^1
  = \frac{2(1 - \cos(m\pi))}{m\pi}
  = \frac{2(1-(-1)^m)}{m\pi}
\]
Since $B_m = 0$ for even $m$, substituting $m = 2k-1$ and plugging
into the collapsed solution gives the final answer
\[
  u(t,x) = \frac{4}{\pi}\sum_{k=1}^{\infty}
  \frac{1}{2k-1}\,e^{-2\pi^2(2k-1)^2 t}\sin\!\left((2k-1)\pi x\right)
\]
\question[4]{}
We now consider the same PDE over the same domain except the
Dirichlet and Neumann conditions are flipped:

We solve the following PDE on the rectangle $-1 < x < 1, 0 < y < 1$
\[
  \begin{cases}
    u_t = 2\Delta u\\
    u(t,x,0) = u(t, x, 1) = 0 & \forall y \\
    u_n(t,-1,y) = u_n(t,1,x) = 0 \\
  \end{cases}
\]
With initial distribution:
\[
  u(0,x,y) =
  \begin{cases}
    -1 & x < 0 \\
    1 & x > 0
  \end{cases}
\]

We skip to the possible solutions of the seperated ODE:

\[
  p = \{A\cos\!\left(\sqrt{\tfrac{\mu}{2}}\,x\right),
  B\sin\!\left(\sqrt{\tfrac{\mu}{2}}\,x\right)\}
\]

\[
  q = \{C\cos\!\left(\sqrt{\tfrac{\lambda-\mu}{2}}\,y\right) ,
  D\sin\!\left(\sqrt{\tfrac{\lambda-\mu}{2}}\,y\right)\}
\]

Now the normal vectors of the short edges are given by $(1,0)$ and
$(-1,0)$, so the Neumann conditions are:
\begin{align*}
  &\nabla u \cdot (-1,0)\big|_{x=-1} = -u_x(t,-1,y) = -e^{-\lambda
  t}p'(-1)q(y) = 0 \implies p'(-1) = 0\\
  &\nabla u \cdot (1,0)\big|_{x=1} = u_x(t,1,y) = e^{-\lambda
  t}p'(1)q(y) = 0 \implies p'(1) = 0
\end{align*}

Applying these to $p' =
-A\sqrt{\tfrac{\mu}{2}}\sin\!\left(\sqrt{\tfrac{\mu}{2}}\,x\right) +
B\sqrt{\tfrac{\mu}{2}}\cos\!\left(\sqrt{\tfrac{\mu}{2}}\,x\right)$:
\begin{align*}
  p'(1)  &=
  -A\sqrt{\tfrac{\mu}{2}}\sin\!\left(\sqrt{\tfrac{\mu}{2}}\right) +
  B\sqrt{\tfrac{\mu}{2}}\cos\!\left(\sqrt{\tfrac{\mu}{2}}\right) = 0 \\
  p'(-1) &=
  \phantom{-}A\sqrt{\tfrac{\mu}{2}}\sin\!\left(\sqrt{\tfrac{\mu}{2}}\right)
  + B\sqrt{\tfrac{\mu}{2}}\cos\!\left(\sqrt{\tfrac{\mu}{2}}\right) = 0
\end{align*}
Adding gives
$2B\sqrt{\tfrac{\mu}{2}}\cos\!\left(\sqrt{\tfrac{\mu}{2}}\right) =
0$, so either $B=0$ or
$\sqrt{\tfrac{\mu}{2}} = \tfrac{(2m-1)\pi}{2}$, giving $p_m =
\sin\!\left(\tfrac{(2m-1)\pi}{2}x\right)$. Subtracting gives
$2A\sqrt{\tfrac{\mu}{2}}\sin\!\left(\sqrt{\tfrac{\mu}{2}}\right) =
0$, so either $A=0$ or
$\sqrt{\tfrac{\mu}{2}} = m\pi$, giving $p_m = \cos(m\pi x)$ for $m =
0,1,2,\ldots$

We now apply the Dirichlet conditions on the long edges:
\begin{align*}
  q(0) = C\cos\!\left(0\right) +
  D\sin\!\left(0\right) &=0 \\
  q(1) = C\cos\!\left(\sqrt{\frac{\lambda - \mu}{2}}\right)
  +D\sin\!\left(\sqrt{\frac{\lambda - \mu}{2}}\right) &=0 \\
\end{align*}

From the first equation we clearly have that $A = 0$, therefore,
using the second equation we have that
$$
\sqrt{\frac{\lambda-\mu}{2}} = n\pi
$$
For $n = 1,2,\ldots$, giving eigenfunctions $q_n(y) = \sin(n\pi
y)$. We can solve for $\mu$ and $\lambda$: $\lambda = \mu + 2n^2\pi^2$: for
the cosine $p$-family
$\mu = 2m^2\pi^2$, and for the sine $p$-family $\mu =
\tfrac{(2m-1)^2\pi^2}{2}$. The general solution is:
\[
  u(t,x,y) =
  \sum_{m=0}^{\infty}\sum_{n=1}^{\infty}
  e^{-2\pi^2(m^2+n^2)t}
  A_{mn}\cos(m\pi x)\sin(n\pi y)
  +
  \sum_{m=1}^{\infty}\sum_{n=1}^{\infty}
  e^{-\left(\frac{(2m-1)^2\pi^2}{2} + 2n^2\pi^2\right)t}
  B_{mn}\sin\!\left(\tfrac{(2m-1)\pi}{2}x\right)\sin(n\pi y)
\]

At $t = 0$

\[
  f(x) =
  \sum_{m=0}^{\infty}\sum_{n=1}^{\infty}
  A_{mn}\cos(m\pi x)\sin(n\pi y)
  +
  \sum_{m=1}^{\infty}\sum_{n=1}^{\infty}
  B_{mn}\sin\!\left(\tfrac{(2m-1)\pi}{2}x\right)\sin(n\pi y)
\]
Since $f$ is odd in $x$ and $\cos(m\pi x)$ is even, orthogonality forces
$A_{mn}=0$ for all $m,n$.
\[
  f(x) =
  \sum_{m=1}^{\infty}\sum_{n=1}^{\infty}
  B_{mn}\sin\!\left(\tfrac{(2m-1)\pi}{2}x\right)\sin(n\pi y)
\]

By a similar exploitation of orthogonality as the prior question, we have that:

\[
  B_{m,n} = 2\left\langle
  f,\,\sin\!\left(\tfrac{(2m-1)\pi}{2}x\right)\sin(n\pi y)\right\rangle
  = 2\!\int_{-1}^{1}\!\int_0^1 f(x)
  \sin\!\left(\tfrac{(2m-1)\pi}{2}x\right)\sin(n\pi y)\,dy\,dx
\]
Since $f$ is independent of $y$ the integrals separate:
\[
  B_{m,n} = 2\left(\int_{-1}^1 f(x)
  \sin\!\left(\tfrac{(2m-1)\pi}{2}x\right)dx\right)
  \left(\int_0^1\sin(n\pi y)\,dy\right)
  = \frac{8(1-(-1)^n)}{(2m-1)n\pi^2}
\]
Since the exponential factors in $m$ and $n$
independently, the final solution is:
\[
  u(t,x,y) =
  \left(\frac{4}{\pi}\sum_{m=1}^{\infty}
    \frac{e^{-\frac{(2m-1)^2\pi^2}{2}t}}{2m-1}
  \sin\!\left(\tfrac{(2m-1)\pi}{2}x\right)\right)
  \!\left(\frac{4}{\pi}\sum_{n=1}^{\infty}
    \frac{e^{-2(2n-1)^2\pi^2 t}}{2n-1}
  \sin\!\left((2n-1)\pi y\right)\right)
\]

\end{document}
